\section{Introduction}

Electronic dance music often combines a strong rhythmic pulse with a wide range of emotional character. Our interest in this project developed from listening to songs that we perceive as sad despite their energetic production. This raises a question that can be approached through statistical learning: which musical characteristics help explain the emotional mood of EDM?

We investigate relationships between measured acoustic features and supplied valence and arousal scores. Valence describes emotional positivity, while arousal describes activation. Separating these responses makes it possible to study negative mood alongside energetic sound. It also avoids defining emotional arousal directly from an audio measurement such as loudness or tempo.

The present report analyzes 92 provisional dance-style candidates selected from MERGE Audio Complete. We extract interpretable measurements of amplitude, rhythm, spectral brightness, spectral change, and bass-band content, then fit separate ridge regressions for valence and arousal. Grouped evaluation keeps each complete artist-credit group within a single fold. An exploratory extension measures temporal changes within the excerpts and connects those measurements to the later study of buildups, drops, and breakdowns.

The initial results show a modest improvement over a training-mean baseline for arousal, with the clearest contribution from spectral brightness and spectral change. The same seven-feature model does not improve valence prediction. Coarse temporal descriptors produce small, split-dependent changes in performance. These findings motivate a more targeted examination of musical structure and additional characteristics, especially those related to harmony and melody. The final EDM review and the two personal song case studies remain part of the next phase of the project.

Section 2 reviews the relevant literature. Section 3 defines the model, documents the data selection and implementation, and presents the evaluated results. The conclusion summarizes what these results support and the further evidence needed to address the case-study questions.
